\subsection{对比实验}
\label{sec:comparison}

将本文方法与四类基线在完全相同的 464 个场景、同一真值与同一评价脚本下比较：（1）Argmax，原始直方图直接取峰（无 Coates 校正与匹配滤波）；（2）Raw+MF，原始直方图经匹配滤波后取峰；（3）Coates+MF，即消融实验中的 L0；（4）CASPI~\cite{Lee2023}，使用官方实现与原论文参数；（5）本文方法（L6，$K=5$）。结果见表~\ref{tab:comp230}与表~\ref{tab:comp46}。

\begin{table}[t]
\centering
\caption{230×310 分辨率下与基线方法的定量对比（464 个场景全集；上半部分为 MAE，单位 m；下半部分为 inl@1\%，单位 \%；列组标题为全像幅平均检测 SNR，mean$\pm\mathrm{std}$，单位 dB）}
\label{tab:comp230}\scriptsize\setlength{\tabcolsep}{6.5pt}
\begin{tabular}{l cc cc cc}
\toprule
 & \multicolumn{2}{c}{$-5.7 \pm 8.4$~dB} & \multicolumn{2}{c}{$-8.4 \pm 8.6$~dB} & \multicolumn{2}{c}{$-15.2 \pm 9.0$~dB} \\
\cmidrule(lr){2-3} \cmidrule(lr){4-5} \cmidrule(lr){6-7}
Method & MAE & inl@1\% & MAE & inl@1\% & MAE & inl@1\% \\
\midrule
Argmax    & 1.645 & 18.28 & 1.724 & 14.85 & 1.976 & 7.19 \\
MF        & 0.404 & 39.56 & 0.553 & 33.14 & 1.048 & 18.08 \\
Coates+MF & 0.665 & 37.04 & 0.814 & 30.94 & 1.694 & 12.48 \\
CASPI     & 0.207 & 38.07 & 0.282 & 35.02 & 0.648 & 25.43 \\
Ours      & \textbf{0.184} & \textbf{40.12} & \textbf{0.207} & \textbf{38.22} & \textbf{0.427} & \textbf{27.46} \\
\bottomrule
\end{tabular}
\end{table}

\begin{table}[t]
\centering
\caption{46×62 分辨率下与基线方法的定量对比（464 个场景全集；上半部分为 MAE，单位 m；下半部分为 inl@1\%，单位 \%；列组标题为全像幅平均检测 SNR，mean$\pm\mathrm{std}$，单位 dB）}
\label{tab:comp46}\scriptsize\setlength{\tabcolsep}{6.5pt}
\begin{tabular}{l cc cc cc}
\toprule
 & \multicolumn{2}{c}{$-5.9 \pm 8.7$~dB} & \multicolumn{2}{c}{$-8.6 \pm 8.9$~dB} & \multicolumn{2}{c}{$-15.4 \pm 9.4$~dB} \\
\cmidrule(lr){2-3} \cmidrule(lr){4-5} \cmidrule(lr){6-7}
Method & MAE & inl@1\% & MAE & inl@1\% & MAE & inl@1\% \\
\midrule
Argmax    & 1.645 & 18.25 & 1.728 & 14.76 & 1.980 & 7.17 \\
MF        & 0.415 & 43.14 & 0.559 & 36.26 & 1.073 & 20.20 \\
Coates+MF & 0.716 & 39.84 & 0.860 & 33.48 & 1.727 & 13.23 \\
CASPI     & \textbf{0.097} & 53.39 & \textbf{0.140} & 51.65 & \textbf{0.375} & 43.63 \\
Ours      & \textbf{0.058} & \textbf{64.28} & \textbf{0.074} & \textbf{63.95} & \textbf{0.186} & \textbf{57.60} \\
\bottomrule
\end{tabular}
\end{table}

对比给出三点观察。首先，传统预处理的选择比直觉敏感：Coates 校正作为反卷积类操作，在光子饥饿条件下将稀疏计数中的噪声放大为高频伪峰，使 Coates+MF 系统性劣于不校正的 Raw+MF（MAE 劣化 $47\%\sim73\%$），匹配滤波在此承担的实为抑制该校正副作用的角色；该差异在 argmax 级同样存在（不校正直方图直接取峰优于 Coates 后取峰 2 倍以上），说明这是校正本身的统计性质而非某一实现的问题。其次，CASPI 在 46×62 上是强基线（较单像素基线改善 $55\%\sim69\%$），但在 230×310 上相对自身退化、与本文方法的差距拉大至 $40\%\sim51\%$：其非局部校正依赖像素间统计相似性，原生分辨率下单像素光子数下降一个量级后相似性证据随之减弱，而本文方法因单目先验提供了与光子统计无关的几何约束，受分辨率与照度的影响显著更小。最后，在全部 12 个工作点中本文方法均取得最优 MAE 与 inl@1\%，且尾部指标的领先幅度大于均值指标（230×310 下 inl@1\% 较 CASPI 高 $6\sim14$ 个百分点），与消融实验中融合权重对不可融合邻居的抑制作用一致。效率方面，在相同的 RTX 2080 Ti 上，CASPI 处理单帧需 89.6 s（230×310），本文完整融合流水线（含单目推理）约 55 s，且该耗时同时完成了全部八个消融变体的计算。
